\chapter*{Postscript: The category we haven't mentioned}
\addcontentsline{toc}{chapter}{Postscript}
\markboth{Postscript}{Postscript}

\noindent This book has been about categories inside languages~-- nouns, countability, definiteness, grammaticality~-- and the mechanisms that maintain them. It hasn't said a word about the category \term{language} itself.

That's deliberate, but it shouldn't go unnoticed.

The essentialist temptation is obvious. \mention{English} names a thing; the thing has boundaries; speakers are either inside or outside. Mutual intelligibility is the usual candidate for the gatekeeping test, and it fails in the usual way: Danish and Norwegian are mutually intelligible, Mandarin and Cantonese aren't, yet the former are two languages and the latter one. Weinreich's observation~-- that a language is a dialect with an army and a navy~-- is the nominalist counter-move, and it has the same structure as Haspelmath's position on grammatical categories: the label is a convenience imposed from outside, not a natural boundary discovered within.

The HPC analysis writes itself. A standard language is a property cluster~-- shared phonology, overlapping lexicon, mutual grammaticality judgments, common orthography, institutional recognition, a literary tradition~-- maintained by mechanisms: formal education, broadcasting, publishing norms, iterated interaction within communities, and the gatekeeping of standardisation bodies. Remove the mechanisms and the cluster frays. Dialect continua are the predicted outcome: where institutional maintenance thins, boundaries soften, and the sharp lines on the map turn out to be political artefacts rather than linguistic ones.

The failure modes transfer too. Thin clustering: pidgins share a lexical base but lack the morphosyntactic and phonological depth of the source languages. Fat clustering: \mention{Chinese} lumps Sinitic varieties that share a writing system but diverge on every spoken diagnostic. Negative diagnosis: constructed languages like Esperanto are defined by what they lack (irregular morphology, dialectal variation) rather than by a maintained profile that projects.

The mechanisms are familiar in outline~-- alignment, entrenchment, transmission, gatekeeping~-- but they operate at a different scale, and several of them (media markets, state education policy, migration patterns) are sociolinguistic in ways that the grammatical mechanisms in this book mostly aren't. Working out the details would mean engaging seriously with variationist sociolinguistics, language policy, and historical linguistics in ways the present book doesn't attempt.

That's another book. What matters here is that the framework doesn't flinch at the question. If \mention{English} names a maintained cluster that projects~-- that supports induction about what a new speaker will do, what a new text will look like, what a child acquiring it will converge on~-- then it's a category in the same sense that \mention{noun} is: real at the scale where the maintenance operates, dissolving where it doesn't.

The sustaining, again, is the reality. It just operates at a larger grain.
